\documentclass[11pt]{article}

\usepackage[margin=1in]{geometry}
\usepackage{amsmath,amssymb,amsthm}
\usepackage{hyperref}
\usepackage{microtype}

\title{\textbf{Admissible Mixed Hodge Fixed-Part}}
\author{Inacio F.\ Vasquez \\ Independent Researcher}
\date{}

% --- Theorem Environments ---
\theoremstyle{plain}
\newtheorem{theorem}{Theorem}
\newtheorem{lemma}[theorem]{Lemma}
\newtheorem{corollary}[theorem]{Corollary}

\theoremstyle{remark}
\newtheorem{remark}[theorem]{Remark}

% --- Notation ---
\newcommand{\Q}{\mathbb{Q}}
\newcommand{\mathbbV}{\mathbb{V}}

\begin{document}

\maketitle

\begin{center}
\textbf{Status summary.}
\begin{itemize}
\item Monodromy invariants, fixed-part, and horizontal-subtori statements remain conjectural for admissible VMHS.
\item Current boundaries require separate theorems for the monodromy-invariant sub-VMHS structure, the admissible fixed part, and the identification $T=\mathsf{Flip}(T)$.
\item Lean entries are target references only: \texttt{URF\_Core.lean → Hodge.*}
\end{itemize}
\end{center}

\tableofcontents

\section{Monodromy Invariants}

\begin{conjecture}[Monodromy invariants]\label{lem:hodge-monodromy}
For an admissible $\mathbbV$, flat sections invariant under the Gauss--Manin connection form a sub-VMHS over $\Q$.
\end{conjecture}

\noindent
BOUNDARY: the existence of a Deligne canonical extension with compatible
weight and Hodge filtrations does not by itself prove that invariant flat
sections are closed under those filtrations or satisfy the full sub-VMHS
axioms.  A separate monodromy-invariant sub-VMHS theorem is required.

\noindent
\textbf{Lean target reference:}
\texttt{URF\_Core.lean → Hodge.monodromy\_invariants}

\section{Fixed Part Theorem}

\begin{conjecture}[Fixed part for admissible VMHS]\label{thm:hodge-fixed-part}
Flat sections $H^0(S, \mathbbV_\Q)$ form a rational mixed Hodge substructure.
\end{conjecture}

\noindent
BOUNDARY: admissibility and functoriality of the Gauss--Manin connection do not
by themselves prove that global flat sections inherit compatible weight and
Hodge filtrations, are rational, and satisfy the mixed Hodge substructure
axioms.  A separate fixed-part theorem for admissible VMHS is required.

\noindent
\textbf{Lean target reference:}
\texttt{URF\_Core.lean → Hodge.fixed\_part}

\section{Rationality of Horizontal Subtori}

\begin{conjecture}[Rational horizontal subtori]\label{cor:hodge-rational}
If $T$ is a horizontal real subtorus generated by an admissible normal function, then $T$ is rational: $T = \mathsf{Flip}(T)$.
\end{conjecture}

\noindent
BOUNDARY: the fixed-part theorem identifies a rational mixed Hodge substructure
of flat sections, but it does not by itself identify the original real subtorus
$T$ with its rational envelope $\mathsf{Flip}(T)$.  A separate admissible
rigidity theorem proving this identification is required.

\noindent
\textbf{Lean target reference:}
\texttt{URF\_Core.lean → Hodge.rational\_horizontal\_subtori}

\begin{remark}[Scope]
All results restricted to admissible VMHS.  
\textbf{Frozen module:} Yes.
\end{remark}

\end{document}

